% Conclusion -- drafted by loop0015 item 09 (2026-11-05), in plain English
% (owner's rule of 2026-11-05 05:30). The second subsection reviews every
% entry of paper3/ETH_AMBIGUOUS.MD written up to this section (S1-S8,
% 01a-08xd, and 09a-09d of this item) and lists the ones that affect the
% text as open, by entry number; the owner resolves an entry by editing the
% paper and striking it. Nothing here is a new number: every figure quoted
% is a later section's, with its section named.
\section{Conclusion}
\label{sec:conclusion}
\begingroup\sloppy

OpenAI's family 378 holds a theorem the Lean kernel accepts, and the
theorem refutes the Exponential Time Hypothesis as stated. We read it,
inspected its proof, rebuilt the proof from source and re-checked it
with a second checker, and found no flaw. That is the first thing this
paper says, and it is said without qualification. The second thing is
what the title says, and it is a reading of the statement, not of the
proof. The statement is a guarantee: a machine exists and halts within
a bound. The bound's constants exist and are never computed; where the proof
writes values down, its construction has at least 257 for the
exponent's constant, and the theorem it exhibits has at least 258
there and at least 28,701 in front (Sections~\ref{sec:statement} and
\ref{sec:proof}). The exponent is the manuscript's weaker one, and with
its constant set to one it exceeds the number of variables at every
size from 50 up to a number with 2,391 digits
(Section~\ref{sec:statement}). So at every practical size the verified
guarantee is weaker than the guarantee of trying every assignment. The
manuscript's sharper exponents, which would change that sentence, are
not in the Lean library (Section~\ref{sec:beyond}). What follows from
the verified statement is exactly what the manuscript says follows: the
hypotheses fall, the lower bounds derived from them lose their premise,
and the conclusions of those lower bounds are untouched
(Section~\ref{sec:consequences}). For anyone who runs a solver, nothing
changes (Section~\ref{sec:practitioner}). Two facts are true at
once, and neither diminishes the other. In the limit, the exponent is
smaller than the number of variables by any factor one names; that is
a theorem in our own Lean (Section~\ref{sec:lean}). On every size anyone
will write down, it is larger.

\subsection{What is claimed, and what is not}
\label{sec:conclusion-claims}

The introduction's two lists (Section~\ref{sec:intro-claims}) are the
paper's claims and non-claims, and each later section repeats its own
under ``what this does not claim''. In one sentence each. The verified
statement is a guarantee with uncomputed constants, above $2^n$ at every
practical size with the constants set to one, and it refutes ETH and its
relatives nonetheless. The proof shows no flaw by inspection, builds
from source, and passes two checkers. The sharper exponents are
paper-only. And no running time of any machine, no flaw in anything,
no value of any constant beyond the proof's own, and no refutation of
any derived lower bound's conclusion is asserted. Every number in the
paper has a script under \repofile{paper3/scripts/} that regenerates
it, every quotation has a page and line or a file and line, and the
replay's logs and hashes are under \repofile{paper3/replay/}.

\subsection{The questions this reading left open, by entry}
\label{sec:conclusion-open}

The paper was written under a rule that no question stops the work.
A choice between two sound ways of saying or doing a thing is made and
recorded in \repofile{paper3/ETH_AMBIGUOUS.MD} as one numbered entry.
The entry says what was ambiguous, what was chosen and what the author
may want to change. Then the writing goes on. Sixty-one entries were
written in all: forty-four before the first draft of this section,
four by that draft, eleven by the reading of the whole paper and the
write-up that followed it, and two by the final pass. We reviewed all
sixty-one and sort them here. Twenty-two need no decision. Seven were
resolved by the author in writing: S1, the author line; S3, where the
companion Lean lives, with 07a following it; S4, no width references;
S6, the model counters; S7, the phrase ``for every $50 \le n <
e^{5504}$''; S8, the replay's Lean version. Fifteen are procedural, or
are corrections the text already carries: S5, the per-item time cap of
the writing driver; 01a, the sandbox's core cap; 03b, the manuscript
has 67 pages; 04b, 04d and 04e, three counts corrected from an earlier
audit; 06b, a counter installed for Section~\ref{sec:practitioner};
08xb, which sections the plain-English rewrite covered; 10c, the long
sentences the first reading left, since split; 11c, two audit files
annotated rather than edited; 11d, the later entries listed in the
repository's README, and since here too; 11e, a paragraph for the
repository's instructions file, drafted and not applied; 11f, the
plan's record of what was done; 12a and 12b, the scope of the final
pass and this list's own re-count. The thirty-nine below affect the
text and are open. In each case the text stands as chosen, and the
entry says what would change.

\begin{description}\setlength{\itemsep}{1pt}
\item[The title and its figures.] S2: $2^n$ is typeset with a
  superscript and the ellipsis as three dots; the plain text is in the
  PDF's metadata. 02a: Section~\ref{sec:statement} carries both the
  two-panel Figure~\ref{fig:central} and Figure~\ref{fig:runtime-scale},
  three drawings of one function. 02b: the central figure's peak is
  placed analytically and its legend carries the $+16$. 10b: the
  title's ``is slower'' is reconciled with the paper's non-claim of any
  running time by Sections~\ref{sec:statement} and
  \ref{sec:practitioner}, and by nothing in the title itself.
\item[The statement (Section~\ref{sec:statement}).] 02c: a proof-side
  fact, the constant $256 + A_t$, is quoted in the statement section
  for its value. 02d: the Lean theorem is transcribed as a display, not
  shown verbatim, because the house build cannot typeset its symbols.
  08xd: Section~\ref{sec:statement} says the proof's constant is at
  least 257 and the exhibited theorem's at least 258, two readings of
  two lines, both kept.
\item[Beyond the statement (Section~\ref{sec:beyond}).] 03a: a third
  tag, ``supported'', between verified and paper-only. 03c: what the
  verified exponent would give through Corollary~1.4's argument is not
  worked out. 03d: Corollary~1.2 is stated with the verified exponent
  and tagged supported, not verified.
\item[The proof and the replay (Sections~\ref{sec:proof} and
  \ref{sec:replay}).] 01b: the second checker is \texttt{nanoda\_lib}
  at the commit named, with \texttt{lean4export} at the nearest tag.
  01c: that checker was run with the three axioms permitted and unused
  axioms ignored. 01d: ``replayed'' means the build, the axiom output
  and the second checker, not the comparator run. 04a: thirty-five
  definitions compared, thirty-three identical, two up to a keyword;
  the audit's ``33'' is explained, not kept. 04c: the exhibited
  constants, $A \ge 258$ and $C \ge 28{,}701$, are stated with their
  lines.
\item[What follows (Section~\ref{sec:consequences}).] 05a: the
  manuscript's list of problems is elided from one quotation, by the
  paper's scope. 05b: SETH is refuted by reading the statement at each
  width, and ``SETH implies ETH'' is cited for the record. 05c: the
  hypotheses' sources are copied from the manuscript's reference list
  and not re-read. 05d: two paragraphs are ours and not the
  manuscript's, and say so.
\item[For the practitioner (Section~\ref{sec:practitioner}).] 06a: the
  two guarantee curves share one scale factor, so that they cross where
  the statement puts the crossing. 06c: seven references for the tools
  and the threshold were typed from memory; their DOIs have since been
  checked (10a). 06d: the counters are reported at their defaults,
  including a plateau that looks like a preprocessing budget.
\item[Our Lean (Section~\ref{sec:lean}).] 07b: the axiom audit prints
  two lines with fewer than the three axioms, and ``every line on the
  three'' is read as ``nothing beyond the three''. 07c: the re-diff also
  matches the restatement's body to the theorem's, beyond what was
  asked.
\item[Related work (Section~\ref{sec:related}).] 08a: the audit paper's
  title is typeset with an apostrophe the arXiv record lacks. 08b: four
  small public records were fetched beside the three allowed pages, for
  their dates and commits. 08c: fourteen references carried no checked
  DOI at the time; thirteen have since been found (10a). 08d: ``we
  found no earlier statement'' is bounded to six texts.
\item[Plain English, and this section.] 08xa: the dense passages were
  moved into Appendix~\ref{app:tech}, not cut. 08xc: the glosses are
  repeated in each section and now also in the introduction. 09a: the
  introduction and abstract state the title's third clause in words,
  with no formula beyond $2^n$, by the plain-English rule rather than
  by the item's brief. 09b: the abstract carries two numbers, both from
  Section~\ref{sec:statement}'s script, and no count from the replay.
  09c: this sorting of the entries into resolved, procedural and open
  is the author's to redo. 09d: ``people who checked'' in the stance
  sentence is read as the author and the assistant the acknowledgements
  disclose.
\item[The write-up around the paper.] 10a: thirteen DOIs were added
  from a registry the download rule did not name. 10d: the reading's
  findings that were recorded and not fixed at the time; the two it
  names have since been written and annotated. 11a: the data statement
  takes its repository address and licence from our first paper. 11b:
  the acknowledgements name the second reading as a second session of
  the same assistant. 11g: the references for Lean~4 and Mathlib are
  copied from our first paper's bibliography.
\end{description}

Fourteen of the thirty-nine are bibliographic or typographic: S2, 02b,
02d, 05c, 06c, 08a, 08b, 08c, 11g, and the five spelling, placement and
tagging choices of 02a, 03a, 03d, 04a, 08xa. The rest are choices
about what the text says. Each is a sentence that can be struck
without touching any number.

\subsection{What this does not claim}
\label{sec:conclusion-not}

\begin{itemize}
\item \textbf{Nothing new.} This section states no fact that a
  previous section does not state with its source; the numbers in it
  are Sections~\ref{sec:statement}, \ref{sec:proof} and
  \ref{sec:replay}'s.
\item \textbf{The list of open questions is complete to the day it was
  last re-counted}, 2026-11-05, and includes the entries of the reading
  of the whole paper and of the final pass. Any entry added after that
  day is in the file and not here.
\item \textbf{No question is answered here.} Each entry records a
  choice and what would change; the author's answer, where there is
  one, is in the file's ``Owner resolutions'' and in the text.
\item \textbf{No claim about the working group.} The collection says a
  group of computer scientists will examine the proof
  (Section~\ref{sec:related}); nothing of theirs is on record here, and
  their reading of the proof will carry more weight on the proof than
  our inspection does.
\end{itemize}
\par\endgroup
